\begin{figure}[!htbp]
\centering
% Positions of the copies, in units of 10^5 columns, from the simulation
% described in Section 7.1 (symbols passed one at a time, N = 10^6).
\begin{tikzpicture}[x=1.35cm,y=1cm,>=stealth]
 \draw[->,black!70] (0,0) -- (10.45,0) node[right,font=\footnotesize] {column};
 \foreach \c in {10,6.18035,3.81967,2.36072,1.45904,0.90176,0.55734,
                 0.34448,0.21290,0.13159,0.08137,0.05030,0.03109,0.01923,
                 0.01190,0.00739,0.00457,0.00284,0.00177,0.00111,0.00070,0.00048}
  \draw[black!80,very thin] (\c,-.09) -- (\c,.09);
 \foreach \a/\b in {10/6.18035,6.18035/3.81967,3.81967/2.36072,
                    2.36072/1.45904,1.45904/0.90176,0.90176/0.55734,
                    0.55734/0.34448,0.34448/0.21290}
  \draw[->,queensUpper,semithick] (\a,.12) to[bend right=55] (\b,.12);
 \node[font=\footnotesize,text=queensUpper,above] at (8.1,.95) {reads};
 \foreach \c/\k/\v in {10/1/{1\,000\,000},6.18035/2/{618\,035},
                       3.81967/3/{381\,967},2.36072/4/{236\,072}}
  \node[below=3pt,align=center,font=\footnotesize] at (\c,0)
   {copy~\k\\$\v$};
 \node[below=3pt,font=\footnotesize] at (0,0) {$0$};
\end{tikzpicture}
\caption{Where the copies work when the outermost copy reaches column
$N=10^6$, with symbols passed one at a time. Each tick marks the next
column of one copy. Each arc leads from where a copy writes to where
it reads, which is where the next copy writes. The ticks crowd towards
the left and end with the $22$nd copy, before column $48$.}
\label{fig:recursive-generation}
\end{figure}
